\documentclass[12pt,a4paper]{article}

% Fonts — LuaLaTeX
\usepackage{fontspec}
\usepackage{unicode-math}
\setmainfont{Linux Libertine O}
\setmathfont{Latin Modern Math}
\newfontfamily\igprimfont{Everson Mono}
\newfontfamily\shavfont[Scale=1.0]{Trabajo}
\setmonofont[Scale=0.85]{DejaVu Sans Mono}
\newcommand{\igfont}{\shavfont}

% Shavian Unicode block (U+10450–U+1047F) → Trabajo automatically
% Works in text mode, inside \text{} in math, and in table cells.
\usepackage{newunicodechar}
\newunicodechar{𐑐}{\ifmmode\text{{\shavfont 𐑐}}\else{{\shavfont 𐑐}}\fi}
\newunicodechar{𐑑}{\ifmmode\text{{\shavfont 𐑑}}\else{{\shavfont 𐑑}}\fi}
\newunicodechar{𐑒}{\ifmmode\text{{\shavfont 𐑒}}\else{{\shavfont 𐑒}}\fi}
\newunicodechar{𐑓}{\ifmmode\text{{\shavfont 𐑓}}\else{{\shavfont 𐑓}}\fi}
\newunicodechar{𐑔}{\ifmmode\text{{\shavfont 𐑔}}\else{{\shavfont 𐑔}}\fi}
\newunicodechar{𐑕}{\ifmmode\text{{\shavfont 𐑕}}\else{{\shavfont 𐑕}}\fi}
\newunicodechar{𐑖}{\ifmmode\text{{\shavfont 𐑖}}\else{{\shavfont 𐑖}}\fi}
\newunicodechar{𐑗}{\ifmmode\text{{\shavfont 𐑗}}\else{{\shavfont 𐑗}}\fi}
\newunicodechar{𐑘}{\ifmmode\text{{\shavfont 𐑘}}\else{{\shavfont 𐑘}}\fi}
\newunicodechar{𐑙}{\ifmmode\text{{\shavfont 𐑙}}\else{{\shavfont 𐑙}}\fi}
\newunicodechar{𐑚}{\ifmmode\text{{\shavfont 𐑚}}\else{{\shavfont 𐑚}}\fi}
\newunicodechar{𐑛}{\ifmmode\text{{\shavfont 𐑛}}\else{{\shavfont 𐑛}}\fi}
\newunicodechar{𐑜}{\ifmmode\text{{\shavfont 𐑜}}\else{{\shavfont 𐑜}}\fi}
\newunicodechar{𐑝}{\ifmmode\text{{\shavfont 𐑝}}\else{{\shavfont 𐑝}}\fi}
\newunicodechar{𐑞}{\ifmmode\text{{\shavfont 𐑞}}\else{{\shavfont 𐑞}}\fi}
\newunicodechar{𐑟}{\ifmmode\text{{\shavfont 𐑟}}\else{{\shavfont 𐑟}}\fi}
\newunicodechar{𐑠}{\ifmmode\text{{\shavfont 𐑠}}\else{{\shavfont 𐑠}}\fi}
\newunicodechar{𐑡}{\ifmmode\text{{\shavfont 𐑡}}\else{{\shavfont 𐑡}}\fi}
\newunicodechar{𐑢}{\ifmmode\text{{\shavfont 𐑢}}\else{{\shavfont 𐑢}}\fi}
\newunicodechar{𐑣}{\ifmmode\text{{\shavfont 𐑣}}\else{{\shavfont 𐑣}}\fi}
\newunicodechar{𐑤}{\ifmmode\text{{\shavfont 𐑤}}\else{{\shavfont 𐑤}}\fi}
\newunicodechar{𐑥}{\ifmmode\text{{\shavfont 𐑥}}\else{{\shavfont 𐑥}}\fi}
\newunicodechar{𐑦}{\ifmmode\text{{\shavfont 𐑦}}\else{{\shavfont 𐑦}}\fi}
\newunicodechar{𐑧}{\ifmmode\text{{\shavfont 𐑧}}\else{{\shavfont 𐑧}}\fi}
\newunicodechar{𐑨}{\ifmmode\text{{\shavfont 𐑨}}\else{{\shavfont 𐑨}}\fi}
\newunicodechar{𐑩}{\ifmmode\text{{\shavfont 𐑩}}\else{{\shavfont 𐑩}}\fi}
\newunicodechar{𐑪}{\ifmmode\text{{\shavfont 𐑪}}\else{{\shavfont 𐑪}}\fi}
\newunicodechar{𐑫}{\ifmmode\text{{\shavfont 𐑫}}\else{{\shavfont 𐑫}}\fi}
\newunicodechar{𐑬}{\ifmmode\text{{\shavfont 𐑬}}\else{{\shavfont 𐑬}}\fi}
\newunicodechar{𐑭}{\ifmmode\text{{\shavfont 𐑭}}\else{{\shavfont 𐑭}}\fi}
\newunicodechar{𐑮}{\ifmmode\text{{\shavfont 𐑮}}\else{{\shavfont 𐑮}}\fi}
\newunicodechar{𐑯}{\ifmmode\text{{\shavfont 𐑯}}\else{{\shavfont 𐑯}}\fi}
\newunicodechar{𐑰}{\ifmmode\text{{\shavfont 𐑰}}\else{{\shavfont 𐑰}}\fi}
\newunicodechar{𐑱}{\ifmmode\text{{\shavfont 𐑱}}\else{{\shavfont 𐑱}}\fi}
\newunicodechar{𐑲}{\ifmmode\text{{\shavfont 𐑲}}\else{{\shavfont 𐑲}}\fi}
\newunicodechar{𐑳}{\ifmmode\text{{\shavfont 𐑳}}\else{{\shavfont 𐑳}}\fi}
\newunicodechar{𐑴}{\ifmmode\text{{\shavfont 𐑴}}\else{{\shavfont 𐑴}}\fi}
\newunicodechar{𐑵}{\ifmmode\text{{\shavfont 𐑵}}\else{{\shavfont 𐑵}}\fi}
\newunicodechar{𐑶}{\ifmmode\text{{\shavfont 𐑶}}\else{{\shavfont 𐑶}}\fi}
\newunicodechar{𐑷}{\ifmmode\text{{\shavfont 𐑷}}\else{{\shavfont 𐑷}}\fi}
\newunicodechar{𐑸}{\ifmmode\text{{\shavfont 𐑸}}\else{{\shavfont 𐑸}}\fi}
\newunicodechar{𐑹}{\ifmmode\text{{\shavfont 𐑹}}\else{{\shavfont 𐑹}}\fi}
\newunicodechar{𐑺}{\ifmmode\text{{\shavfont 𐑺}}\else{{\shavfont 𐑺}}\fi}
\newunicodechar{𐑻}{\ifmmode\text{{\shavfont 𐑻}}\else{{\shavfont 𐑻}}\fi}
\newunicodechar{𐑼}{\ifmmode\text{{\shavfont 𐑼}}\else{{\shavfont 𐑼}}\fi}
\newunicodechar{𐑽}{\ifmmode\text{{\shavfont 𐑽}}\else{{\shavfont 𐑽}}\fi}
\newunicodechar{𐑾}{\ifmmode\text{{\shavfont 𐑾}}\else{{\shavfont 𐑾}}\fi}
\newunicodechar{𐑿}{\ifmmode\text{{\shavfont 𐑿}}\else{{\shavfont 𐑿}}\fi}

% Page layout
\usepackage[top=1in, bottom=1in, left=1in, right=1in]{geometry}
\usepackage{microtype}
\usepackage{parskip}

% Language
\usepackage[english]{babel}

% Headers
\usepackage{fancyhdr}
\pagestyle{fancy}
\fancyhf{}
\fancyhead[L]{\leftmark}
\fancyhead[R]{\thepage}
\fancyfoot[C]{\thepage}
\renewcommand{\headrulewidth}{0.4pt}
\renewcommand{\footrulewidth}{0pt}

% Section styling — fully unnumbered; pipeline never auto-numbers
\usepackage{titlesec}
\setcounter{secnumdepth}{0}
\titleformat{\section}{\Large\bfseries}{}{0em}{}
\titleformat{\subsection}{\large\bfseries}{}{0em}{}
\titleformat{\subsubsection}{\normalsize\bfseries}{}{0em}{}

% Essential packages
\usepackage{graphicx}
\graphicspath{{/home/mrnob0dy666/imsgct/red-hot_rebis/docs/}{/home/mrnob0dy666/imsgct/red-hot_rebis/docs/images/}}
\setkeys{Gin}{width=0.48\linewidth,height=0.32\textheight,keepaspectratio}
\usepackage[table]{xcolor}
\usepackage{booktabs}
\usepackage{array}
\usepackage{tabularx}
\usepackage{longtable}
\usepackage{amsmath}
% unicode-math subsumes amssymb — do not load amssymb alongside it
\usepackage{float}
\usepackage[normalem]{ulem}
\renewcommand{\arraystretch}{1.3}

% Cross-references
\usepackage{hyperref}
\usepackage{cleveref}
\hypersetup{colorlinks=true, linkcolor=blue, urlcolor=cyan, citecolor=teal}

% Custom commands
\newcommand{\shav}[1]{{\shavfont #1}}
\newcommand{\heb}[1]{{\igprimfont #1}}
\newcommand{\tupleaddr}[1]{$\langle #1 \rangle$}

% Shorthand primitive commands — redefined for IG document use
\renewcommand{\H}{Ħ}
\newcommand{\K}{Ç}
\newcommand{\G}{Γ}
\newcommand{\g}{ɢ}
\newcommand{\Th}{Þ}
\newcommand{\D}{Ð}
\newcommand{\R}{Ř}
\newcommand{\f}{ƒ}

% Imscription tuple box
\usepackage{tcolorbox}
\tcbuselibrary{skins,breakable}
\newtcolorbox{imscriptionbox}{enhanced, colback=blue!5, colframe=blue!75!black, fontupper=\large, center upper, boxrule=1pt, arc=4pt}

% PDF metadata
\hypersetup{
  pdftitle={Best Small Molecules},
  pdfkeywords={Imscribing Grammar},
  pdfauthor={Lando Mills},
}

\title{Best Small Molecules}
\date{\today}
\author{Lando Mills}

\begin{document}
\maketitle
\hypertarget{best-small-molecules-designed-in-red-hot-rebis}{%
\section{Best Small Molecules Designed in Red-Hot
Rebis}\label{best-small-molecules-designed-in-red-hot-rebis}}

Author: Lando⊗⊙perator

A curated set of the strongest small molecules \textbf{designed} inside
this repository. These are novel structural designs, not catalogued
natural products and not proteins, genes, or biologics. Each entry gives
the structure, what it targets, its 12-primitive structural type, its
imscription tier, and why it stands out. Every SMILES below is taken
verbatim from the repository sources; nothing is invented.

The designs fall into three families: crystal-navigated autocatalytic
tryptamines, a Frobenius-coupled selective chemotherapeutic, and de-novo
ligands imscribed straight from enzyme active sites.

\begin{center}\rule{0.5\linewidth}{0.5pt}\end{center}

\hypertarget{autocatalytic-tryptamines}{%
\subsection{Autocatalytic tryptamines}\label{autocatalytic-tryptamines}}

These are the repository's headline molecules: the first small molecules
imscribed at (O\_\infty) self-modeling criticality. They were not found
by screening. The crystal navigator located an empty (\odot)-critical
neighborhood and computed the single-primitive promotion 𐑮→⊙
(complex-critical to self-modeling critical) needed to inhabit it, then
read off the molecule that satisfies it.

\hypertarget{nitro-bufotenin-dmt-}{%
\subsubsection{1. 5-Nitro-Bufotenin
(DMT-⊙)}\label{nitro-bufotenin-dmt-}}

\begin{itemize}
\tightlist
\item
  \textbf{SMILES:}
  \texttt{CN(C)CCC1=CNC2=CC=C(O)C({[}N+{]}(=O){[}O-{]})=C12}
\item
  \textbf{Structural type:} \texttt{⟨𐑨𐑸𐑩𐑿𐑐𐑘𐑔𐑵⊙𐑫𐑕𐑭⟩}, criticality 𐑮
  promoted to ⊙
\item
  \textbf{Tier:} (O\_\infty) (self-modeling, autocatalytic)
\item
  \textbf{What it does:} carries all three Dialetheia tokens at once;
  EVALT (phenol) plus EVALF (dimethylamine) plus ENGAGR (nitro). With
  the triad complete, the molecule is already (\odot)-critical with no
  further modification, which is what licenses template-directed
  autocatalytic synthesis: the molecule is its own template.
\item
  \textbf{Why it matters:} the first imscribed small molecule that is
  structurally self-modeling. It was reached by a single clean primitive
  promotion from DMT, so the design and its provenance are both
  univocal. Registered in \texttt{imas/compound\_catalog.py} as the
  DMT-⊙ entry.
\end{itemize}

\hypertarget{nitro-serotonin}{%
\subsubsection{2. 5-Nitro-Serotonin}\label{nitro-serotonin}}

\begin{itemize}
\tightlist
\item
  \textbf{SMILES:}
  \texttt{NCCC1=CNC2=CC=C(O)C({[}N+{]}(=O){[}O-{]})=C12}
\item
  \textbf{Tier:} (O\_\infty) (self-modeling, autocatalytic)
\item
  \textbf{What it does:} the same 𐑮→⊙ promotion applied to serotonin
  rather than DMT, confirming the autocatalytic motif is a transferable
  structural operation and not a one-off coincidence of DMT.
\item
  \textbf{Why it matters:} demonstrates the ⊙-promotion as a reusable
  design move across the tryptamine family. Registered alongside DMT-⊙
  in \texttt{imas/compound\_catalog.py}.
\end{itemize}

\begin{center}\rule{0.5\linewidth}{0.5pt}\end{center}

\hypertarget{frobenius-coupled-selective-chemotherapeutic}{%
\subsection{Frobenius-coupled selective
chemotherapeutic}\label{frobenius-coupled-selective-chemotherapeutic}}

\hypertarget{frobenius-coupled-chemotherapeutic-dimer}{%
\subsubsection{3. Frobenius-Coupled Chemotherapeutic
dimer}\label{frobenius-coupled-chemotherapeutic-dimer}}

\begin{itemize}
\tightlist
\item
  \textbf{Structural type:} \texttt{⟨𐑦𐑶𐑾𐑹𐑐𐑧𐑲𐑝⊙𐑫𐑳𐑭⟩}
\item
  \textbf{Tier:} (O\_\infty)
\item
  \textbf{What it does:} a tether-coupled two-domain molecule whose
  cytotoxic payload stays masked while the host cell satisfies the
  Frobenius check μ∘δ=id. On healthy tissue the payload never exposes,
  so healthy-cell cytotoxicity is exactly zero. Where malignant
  asymmetry breaks the check (μ∘δ≠id), the tether tension crosses
  threshold, the payload releases, and the molecule becomes cytotoxic.
\item
  \textbf{Why it matters:} the selectivity gate is structural, not
  pharmacological. In the bundled simulation the molecule shows zero
  activity on healthy cells and on low-asymmetry tumors, then full
  payload release and a sharp selectivity jump once mutation strength
  passes the asymmetry threshold, with a selectivity ratio in the
  four-figure range over healthy tissue. Source:
  \texttt{therapeutics/frobenius\_chemotherapeutic.py}, results in
  \texttt{therapeutics/frobenius\_chemo\_results.json}.
\end{itemize}

\begin{center}\rule{0.5\linewidth}{0.5pt}\end{center}

\hypertarget{de-novo-enzyme-ligands}{%
\subsection{De-novo enzyme ligands}\label{de-novo-enzyme-ligands}}

These are imscribed directly from enzyme active-site structural types by
the reverse-ligand pipeline
(\texttt{rhr\_p4rky/ligand\_from\_active\_site.py}); no template
libraries, no known-inhibitor seeds, no docking scores. The four below
are the genuinely novel designs from that run. As a validation check,
the same pipeline independently re-derives ethanol (the real
alcohol-dehydrogenase substrate) and biuret (a known urease inhibitor)
from structure alone, and every candidate it produces is
Lipinski-compliant. Source analysis:
\texttt{ig-docs/bevy\_ligand\_analysis/01\_bevy\_ligand\_analysis.md};
ligand set: \texttt{ig-docs/bevy\_final\_ligands.json}.

\hypertarget{cyp2d6-pyridine-ether}{%
\subsubsection{4. CYP2D6 pyridine-ether}\label{cyp2d6-pyridine-ether}}

\begin{itemize}
\tightlist
\item
  \textbf{SMILES:} \texttt{COCCC(C)Cc1cccnc1}
\item
  \textbf{Target:} cytochrome P450 2D6, the enzyme that metabolizes a
  large share of marketed drugs
\item
  \textbf{What it does:} pairs a basic pyridine nitrogen with an
  aromatic ring at the spacing the 2D6 σ-bond hydroxylation
  pharmacophore wants; drug-like and Lipinski-compliant.
\end{itemize}

\hypertarget{carbonic-anhydrase-ii-zinc-coordinator}{%
\subsubsection{5. Carbonic anhydrase II
zinc-coordinator}\label{carbonic-anhydrase-ii-zinc-coordinator}}

\begin{itemize}
\tightlist
\item
  \textbf{SMILES:} \texttt{O=Cc1cccc(C(=O)O)c1}
\item
  \textbf{Target:} carbonic anhydrase II, a zinc metalloenzyme
\item
  \textbf{What it does:} presents a carboxylic acid plus aldehyde for
  coordinating the Zn²⁺-activated hydroxide; a de-novo analog of the
  classical CA inhibitor pharmacophore reached without that prior.
\end{itemize}

\hypertarget{lysozyme-bis-epoxide-strain-release-inhibitor}{%
\subsubsection{6. Lysozyme bis-epoxide strain-release
inhibitor}\label{lysozyme-bis-epoxide-strain-release-inhibitor}}

\begin{itemize}
\tightlist
\item
  \textbf{SMILES:} \texttt{c1cc(C2CO2)cc(C2CC2c2ccc(C3CO3)cc2)c1}
\item
  \textbf{Target:} lysozyme, a glycoside hydrolase
\item
  \textbf{What it does:} the epoxides mimic the flattened oxocarbenium
  transition state of glycoside hydrolysis, and the bis-electrophile
  design offers two handles for covalent capture in the cleft.
\end{itemize}

\hypertarget{petase-amide-pet-mimic}{%
\subsubsection{7. PETase amide PET-mimic}\label{petase-amide-pet-mimic}}

\begin{itemize}
\tightlist
\item
  \textbf{SMILES:} \texttt{NC(=O)c1ccc(NC(=O)c2cccc(C=O)c2)cc1}
\item
  \textbf{Target:} PETase, the engineered plastic-degrading enzyme
\item
  \textbf{What it does:} a de-novo amide ligand shaped to the
  PET-binding groove; the standout environmental-remediation candidate
  of the set.
\end{itemize}

\begin{center}\rule{0.5\linewidth}{0.5pt}\end{center}

\hypertarget{at-a-glance}{%
\subsection{At a glance}\label{at-a-glance}}

\begin{longtable}[]{@{}lllll@{}}
\toprule
\begin{minipage}[b]{0.06\columnwidth}\raggedright
\#\strut
\end{minipage} & \begin{minipage}[b]{0.20\columnwidth}\raggedright
Molecule\strut
\end{minipage} & \begin{minipage}[b]{0.16\columnwidth}\raggedright
Family\strut
\end{minipage} & \begin{minipage}[b]{0.12\columnwidth}\raggedright
Tier\strut
\end{minipage} & \begin{minipage}[b]{0.31\columnwidth}\raggedright
Target / role\strut
\end{minipage}\tabularnewline
\midrule
\endhead
\begin{minipage}[t]{0.06\columnwidth}\raggedright
1\strut
\end{minipage} & \begin{minipage}[t]{0.20\columnwidth}\raggedright
5-Nitro-Bufotenin (DMT-⊙)\strut
\end{minipage} & \begin{minipage}[t]{0.16\columnwidth}\raggedright
autocatalytic tryptamine\strut
\end{minipage} & \begin{minipage}[t]{0.12\columnwidth}\raggedright
(O\_\infty)\strut
\end{minipage} & \begin{minipage}[t]{0.31\columnwidth}\raggedright
self-modeling autocatalysis\strut
\end{minipage}\tabularnewline
\begin{minipage}[t]{0.06\columnwidth}\raggedright
2\strut
\end{minipage} & \begin{minipage}[t]{0.20\columnwidth}\raggedright
5-Nitro-Serotonin\strut
\end{minipage} & \begin{minipage}[t]{0.16\columnwidth}\raggedright
autocatalytic tryptamine\strut
\end{minipage} & \begin{minipage}[t]{0.12\columnwidth}\raggedright
(O\_\infty)\strut
\end{minipage} & \begin{minipage}[t]{0.31\columnwidth}\raggedright
self-modeling autocatalysis\strut
\end{minipage}\tabularnewline
\begin{minipage}[t]{0.06\columnwidth}\raggedright
3\strut
\end{minipage} & \begin{minipage}[t]{0.20\columnwidth}\raggedright
Frobenius-Coupled dimer\strut
\end{minipage} & \begin{minipage}[t]{0.16\columnwidth}\raggedright
selective cytotoxic\strut
\end{minipage} & \begin{minipage}[t]{0.12\columnwidth}\raggedright
(O\_\infty)\strut
\end{minipage} & \begin{minipage}[t]{0.31\columnwidth}\raggedright
tumor-asymmetry-gated payload\strut
\end{minipage}\tabularnewline
\begin{minipage}[t]{0.06\columnwidth}\raggedright
4\strut
\end{minipage} & \begin{minipage}[t]{0.20\columnwidth}\raggedright
CYP2D6 pyridine-ether\strut
\end{minipage} & \begin{minipage}[t]{0.16\columnwidth}\raggedright
de-novo ligand\strut
\end{minipage} & \begin{minipage}[t]{0.12\columnwidth}\raggedright
O₂\strut
\end{minipage} & \begin{minipage}[t]{0.31\columnwidth}\raggedright
P450 2D6 substrate\strut
\end{minipage}\tabularnewline
\begin{minipage}[t]{0.06\columnwidth}\raggedright
5\strut
\end{minipage} & \begin{minipage}[t]{0.20\columnwidth}\raggedright
CA-II zinc-coordinator\strut
\end{minipage} & \begin{minipage}[t]{0.16\columnwidth}\raggedright
de-novo ligand\strut
\end{minipage} & \begin{minipage}[t]{0.12\columnwidth}\raggedright
O₂\strut
\end{minipage} & \begin{minipage}[t]{0.31\columnwidth}\raggedright
carbonic anhydrase II\strut
\end{minipage}\tabularnewline
\begin{minipage}[t]{0.06\columnwidth}\raggedright
6\strut
\end{minipage} & \begin{minipage}[t]{0.20\columnwidth}\raggedright
Lysozyme bis-epoxide\strut
\end{minipage} & \begin{minipage}[t]{0.16\columnwidth}\raggedright
de-novo ligand\strut
\end{minipage} & \begin{minipage}[t]{0.12\columnwidth}\raggedright
O₂\strut
\end{minipage} & \begin{minipage}[t]{0.31\columnwidth}\raggedright
glycoside hydrolase TS analog\strut
\end{minipage}\tabularnewline
\begin{minipage}[t]{0.06\columnwidth}\raggedright
7\strut
\end{minipage} & \begin{minipage}[t]{0.20\columnwidth}\raggedright
PETase PET-mimic\strut
\end{minipage} & \begin{minipage}[t]{0.16\columnwidth}\raggedright
de-novo ligand\strut
\end{minipage} & \begin{minipage}[t]{0.12\columnwidth}\raggedright
O₂\strut
\end{minipage} & \begin{minipage}[t]{0.31\columnwidth}\raggedright
plastic degradation\strut
\end{minipage}\tabularnewline
\bottomrule
\end{longtable}

The two tryptamines are the crown of the set: imscribed small molecules
that sit at self-modeling criticality by a single, traceable structural
promotion. The Frobenius dimer is the strongest therapeutic design,
turning a structural identity check into a selectivity gate. The four
de-novo ligands show the active-site-to-ligand pipeline producing valid,
drug-like chemistry from first principles.

\end{document}
